

















% (*

\documentclass[12pt]{article}

\title{Lestrade:  A framework for mathematical reasoning with extended examples}

\author{Lavinia Randall Holmes}

\usepackage{amssymb}

\begin{document}

\maketitle

\section{Introduction}

In this essay I will review the concepts behind and demonstrate the use of the Lestrade logical framework and the software which implements it.

There are lots of systems in use which resemble this one in the sense that they use dependent type theory to represent proofs as mathematical objects (and to represent mathematical objects proper).  The systems in use now are quite baroque, with lots of built-in concepts and notation.  Lestrade is minimal in size, and much that would be built in in another system needs to be implemented by the user.  This does mean that the user has more options.

Lestrade is intended to implement a philosophical view of what mathematics is.  The other systems (or some of them) may have similar motivations.  We think the minimal character of Lestrade helps to make this point.

This is a literate programming document.  It can be executed:  the type declarations in the embedded blocks of Lestrade code are actually generated by Lestrade.

\section{General remarks about the framework}

It might be better to go directly to examples below and look back at this section as necessary to understand details of what is happening.


Lestrade is a logical framework of a minimal character, implementing a dependent type theory.

An underlying idea is that mathematical proofs (or evidence for the truth of mathematical statements) are themselves mathematical objects.

A Lestrade text is a series of declarations and definitions:  what the software does is check that they are typed correctly.  Lestrade looks like the type declaration system for a programming language,
and one way the project could be extended is to actually develop the language.

\subsection{The object sorts introduced}

We begin by describing the type system.  General Lestrade types are actually called sorts, because types are a particular variety of sort, as we will see.  There are two kinds of sorts, object sorts
and function sorts.  We begin with the object sorts.

There is a sort {\tt obj} of general mathematical objects.  A theory like ZFC in which all objects are of the same type (in the case of ZFC, sets) would probably have sets implemented as objects of sort {\tt obj}.

There is a sort {\tt prop} of propositions.

There is a sort {\tt type} of  ``type labels''.  A typical object of sort {\tt type} would be {\tt Nat}, a label for the type of natural numbers.  We will see in a moment how it is used and why we might call it a label.

There is, for each proposition $p$, a sort {\tt that} $p$ inhabited by evidence for the truth of $p$.  A proof of $p$ would be of sort {\tt that} $p$ but we do not assume that evidence is proof (evidence might be hypothetical:  when we suppose that $p$ is true, we are not necessarily supposing that we have a proof of it, if we are not constructivists.  Lestrade will implement a constructive view, but it does not commit one to a constructive view.

There is, for each type label $\tau$, a sort {\tt in} $\tau$ inhabited by objects of type $\tau$.  So a natural number might be of sort {\tt in Nat}. We comment briefly that we refer to $\tau$ as a ``type label'' mainly to resist the idea that a type is a (usually infinite) completely given collection.  Though Lestrade is not developed to implement a constructivist philosophy (necessarily) it {\em is\/} developed to implement a philosophy which prefers potential infinity to actual infinity.

These are all of the object sorts of Lestrade.  Note that as we elaborate the object sorts {\tt prop} and {\tt type}, more object sorts become available.

\subsection{Declarations, functions, and function sorts}

In a Lestrade environment, there are constants and parameters.

One can issue a command {\tt declare x tau} to introduce a parameter $x$ of an object sort $\tau$ [it is also possible to introduce parameters with function sorts using this command, but we will restrict ourselves for the moment].

The object type $\tau$ may depend on parameters given earlier in the environment.

One can issue a command {\tt postulate f x1 ... xn tau} to declare a function $f$ which sends parameters $(x_1,\ldots,x_n)$ to an output of object sort $\tau$.  The parameters must appear in the order in which they are declared
(so the type $\tau_i$ of $x_i$ may mention parameters $x_j$ only if $j<i$).  The $x_i$'s may be objects or functions!  For the moment we stipulate that no parameter can appear in any $x_i$ or in $\tau$ other than the $x_i$'s [later, we will see that Lestrade allows some arguments to be left implicit if they can be deduced from context].  {\tt postulate a tau} with no parameters is a permitted variant (declaring an object constant).

One can also issue a command {\tt define f x1 ... xn: t} which will define $f(x_1,\ldots,x_n)$ as $t$ if such a declaration is compatible with the type system.  {\tt define T tau} is a permitted variant.

The reason for describing the basic declaration commands is that this actually specifies what the function sorts of Lestrade are.

A function sort takes the shape $(x_1:\tau_1,\ldots,x_n:\tau_n \Rightarrow \tau)$ where the $x_i$'s are variables bound in the notation, the type $\tau_i$ assigned to $x_i$ will only mention parameters $x_j$ with $j<i$,
and any of the $x_i$'s may be mentioned in $\tau$.  The types $\tau_i$ may be object or function sorts:  $\tau$ must be an object sort.

We describe the conditions under which a term $f(t_1,\ldots,t_n)$ is well-formed and what its sort is if it is well-formed, if $f$ has sort $(x_1:\tau_1,\ldots,x_n:\tau_n \Rightarrow \tau)$.  $t_1$ must have sort
$\tau_1$.  Each $t_i$ with $i>1$ must have sort $\tau_i[t_1/x_1,\ldots t_{i-1}/x_{i-1}]$, the result of replacing each $x_j$ for $j<i$ with $t_j$ in $\tau_j$.  The sort of $f(t_1,\ldots,t_n)$ will then be $\tau[t_1/x_1,\ldots,t_n/x_n]$.

It will, we hope, be possible for the reader to see that this is precisely the range of functions which can be declared with the {\tt postulate} command with the format described above.

Our notion of substitution is aggressive in the sense that substitution of a defined function for a function variable will cause definition expansion.  A command {\tt define f x1 ... xn: t}  will succeed if
$t$ has object sort $\tau$:  this will define a function $f$ of sort $(x_1:\tau_1,\ldots,x_n:\tau_n \Rightarrow \tau)$ which we write $(x_1:\tau_1,\ldots,x_n:\tau_n \Rightarrow t:\tau)$.  If we have an expression
$g(t_1,\ldots,t_n)$ where $g$ is a variable of sort $(x_1:\tau_1,\ldots,x_n:\tau_n \Rightarrow \tau)$, the result of replacing $g$ with $f$ will be the result of replacing each $x_i$ with $t_i$ in $t$.

\subsection{The environment system}

A Lestrade environment, as we have said above, includes constants and includes a list of declared parameters and definitions which may be used in declaration and definition commands as indicated above.

The basic system of environments is actually a list of environments called {\em moves\/}.   When the system is in move $i$, there is a list of moves $0,1,\ldots,i-1,i$.   $i$ is always positive.

If $i<j$, parameters in move $i$ are constants in move $j$.  So when we are in move $i$, our parameters are those declared in move $i$, and parameters declared in move $j<i$ are constants.  When we are in move $i$ and issue the command
{\tt declare x tau}, $x$ is declared with type $\tau$ as a parameter in move $i$.  When we are in move $i$ and issue the command {\tt postulate f x1 ... xn tau}, the arguments $x_1,\ldots,x_n$ must be parameters declared in move $i$, and $f$ is introduced as a parameter in move $i-1$
with sort $(x_1:\tau_1,\ldots,x_n:\tau_n \Rightarrow \tau)$ [and so in move $i$ as a constant].  What the {\tt define x1 ... xn t} command does is similar:  $f$ is introduced as a definition in move $i-1$ with sort
$(x_1:\tau_1,\ldots,x_n:\tau_n \Rightarrow t:\tau)$ [the definition body appears as an annotation in the sort;  a function sort actually has a dummy in place of a definition body in the Lestrade display of its sort].  In any {\tt postulate} or {\tt define} commands issued at move $i$, any definitions appearing in move $i$ must be eliminated by definitional expansion in the declarations placed at move $i-1$. 

The system of environments is active.  There are always at least two moves, move 0 and move 1.  If we are in move $i$ (there are moves with indices from 0 to $i$) and we issue the {\tt open} command,
move $i+1$ is added with no parameter declarations and we are in move $i+1$ [so all previous declarations become constant declarations].  If we are in move $i$ and issue the {\tt clearcurrent} command, all declarations in move $i$ are cleared and we stay in move $i$.  If we are in move $i>1$ and we issue the {\tt close} command, we remove move $i$ and so are in move $i-1$ [{\tt clearcurrent} is needed because we cannot close move 1].  Definitions appearing in move $i$ can be removed safely :  they have been eliminated by definitional expansion in all declarations introduced at move $i-1$ already. 

When a number of environments are opened then closed, definitional expansions can cause considerable increase in size of terms.  This is a technical issue we are working on.

The {\tt close} command is destructive.  There is a system allowing moves to be saved and reopened, which we may discuss later.

The environment system allows something which the discussion above may lead one to believe we do not support.  With the commands introduced so far, it might be hard to see that we can declare function parameters (and you will note that we allow them to be used in function definitions and yet we have only explicitly given a command for declaring object parameters).  The point is that the {\tt postulate} command invoked at move $i$ creates a function parameter 
at move $i-1$.  So one can create a function parameter to be supplied to a further declaration or definition command by using the {\tt open} command followed by suitable declarations of its argument parameters, followed by postulation of the function  and the {\tt close} command.  The objects introduced in move $i-1$ by the {\tt define} command are not general parameters, though there is a method for using them as formal arguments which we are not sure we advocate.

There is a method for declaring function parameters directly, but we will discuss it in the context of examples of use of the system.

\section{What is the difference between a logic and a logical framework?  We implement logic to show you}

In this section, we will implement some logic to show you how the system works and what we are up to.  The opening section is more like a technical reference, and may not really have told you what is happening.

As we do this, we will introduce more technical refinements, and it will be clear why they are needed as we go.

\subsection{Conjunction}

We begin to exhibit Lestrade test.

\begin{verbatim}

begin Lestrade execution

   >>> declare p prop

   p : prop

   {move 1}

   >>> declare q prop

   q : prop

   {move 1}
end Lestrade execution

\end{verbatim}
As will always be the case here, the commands at the prompt {\tt <<<} are emtered by the author, but the responses are from the Lestrade Type Inspector (the implementing software).   Here, we have taken the very simple action of declaring propositions $p$ and $q$.  Notice that these are entered in move 1, as one would expect from the discussion above.

\begin{verbatim}
begin Lestrade execution

   >>> postulate & p q prop

   & : [(p_1 : prop), (q_1 : prop) => 
       (--- : prop)]

   {move 0}
end Lestrade execution
\end{verbatim}

Here we present the demonstration of the operation of conjunction.  The dummy {\tt ---} appears in the sort declaration in the place where the body of the definition would appear in the sort declaration of a defined notion:  it is a signal that {\tt \&} is a primitive notion.

At this point conjunction is nothing but an operation which takes two propositions and returns a proposition.  Further declarations are required to reveal what it is.

\begin{verbatim}
begin Lestrade execution

   >>> declare pp that p

   pp : that p

   {move 1}

   >>> declare qq that q

   qq : that q

   {move 1}
end Lestrade execution
\end{verbatim}

We declare new and different parameters $pp$ and $qq$, evidence for the truth of $p$ and $q$, respectively.

\begin{verbatim}
begin Lestrade execution

   >>> postulate Conj0 p q pp qq that p & q

   Conj0 : [(p_1 : prop), (q_1 : prop), (pp_1 
       : that p_1), (qq_1 : that q_1) => 
       (--- : that p_1 & q_1)]

   {move 0}

   >>> define Conj1 pp qq : Conj0 p q pp \
       qq

   Conj1 : [(.p_1 : prop), (.q_1 : prop), (pp_1 
       : that .p_1), (qq_1 : that .q_1) => 
       ({def} Conj0 (.p_1, .q_1, pp_1, qq_1) : that 
       .p_1 & .q_1)]

   Conj1 : [(.p_1 : prop), (.q_1 : prop), (pp_1 
       : that .p_1), (qq_1 : that .q_1) => 
       (--- : that .p_1 & .q_1)]

   {move 0}

   >>> postulate Conj pp qq that p & q

   Conj : [(.p_1 : prop), (.q_1 : prop), (pp_1 
       : that .p_1), (qq_1 : that .q_1) => 
       (--- : that .p_1 & .q_1)]

   {move 0}
end Lestrade execution

\end{verbatim}

There is a single new idea implemented here, but  we encounter a technical issue.  The function {\tt Conj0} which we declare initially takes propositions $p,q$ and evidence $pp, qq$ for these propositions and returns evidence for $p \wedge q$.  Notice that Lestrade both understands and outputs infix notation {\tt p \& q} though the declaration was in the form {\tt \& p q}.

The operation declared is the rule of conjunction, from $p$ and $q$ deduce $p \wedge q$.

The technical issue is that the arguments $p$ and $q$ appear excess to requirements, since they can be deduced from the types of the subsequent arguments $pp$ and $qq$.  And Lestrade allows us to do this:  there is a mechanism for identifying implicit arguments (found in the types of explicit arguments) and deducing their values.  In complicated cases, the process of deducing values of implicit arguments may fail, and so an explicit form is useful to have available.

Notice that both {\tt Conj1} (defined in terms of {\tt Conj}) and {\tt Conj} (a primitive declared with implicit arguments)
actually have the implicit arguments listed in their type declarations, with an initial period.

There isn't any situation where the implicit arguments cannot be deduced here [this arises when we engage in higher order matching];  we might as well use {\tt Conj} as our basic rule.
\begin{verbatim}

begin Lestrade execution

   >>> declare rr that p & q

   rr : that p & q

   {move 1}

   >>> postulate Simp1 rr that p

   Simp1 : [(.p_1 : prop), (.q_1 : prop), (rr_1 
       : that .p_1 & .q_1) => (--- : that 
       .p_1)]

   {move 0}

   >>> postulate Simp2 rr that q

   Simp2 : [(.p_1 : prop), (.q_1 : prop), (rr_1 
       : that .p_1 & .q_1) => (--- : that 
       .q_1)]

   {move 0}
end Lestrade execution

\end{verbatim}

The functions {\tt Simp1} and {\tt Simp2} allow us to take evidence for {\tt p \& q} and deduce $p$ or $q$ respectively.  Notice the deduction of implicit arguments for these functions.

This actually completes the implementation of conjunction, which is quite a simple logical operation.

We give an example of implementation of a derived rule.

\begin{verbatim}

begin Lestrade execution

   >>> define Selfconj pp : Conj pp pp

   Selfconj : [(.p_1 : prop), (pp_1 
       : that .p_1) => 
       ({def} pp_1 Conj pp_1 : that .p_1 
       & .p_1)]

   Selfconj : [(.p_1 : prop), (pp_1 
       : that .p_1) => (--- : that .p_1 
       & .p_1)]

   {move 0}
end Lestrade execution
\end{verbatim}
and a more complicated example

\begin{verbatim}
begin Lestrade execution

   >>> declare r prop

   r : prop

   {move 1}

   >>> declare ss that (p & q) & r

   ss : that (p & q) & r

   {move 1}

   >>> open

      {move 2}

      >>> define line1 : Simp1 ss

      line1 : that p & q

      {move 1}

      >>> define line2 : Simp1 line1

      line2 : that p

      {move 1}

      >>> define line3 : Simp2 line1

      line3 : that q

      {move 1}

      >>> define line4 : Simp2 ss

      line4 : that r

      {move 1}

      >>> define line5 : Conj line3 line4

      line5 : that q & r

      {move 1}

      >>> define line6 : Conj line2 line5

      line6 : that p & q & r

      {move 1}

      >>> close

   {move 1}

   >>> define Assoctest ss : line6

   Assoctest : [(.p_1 : prop), (.q_1 
       : prop), (.r_1 : prop), (ss_1 
       : that (.p_1 & .q_1) & .r_1) => 
       ({def} Simp1 (Simp1 (ss_1)) Conj 
       Simp2 (Simp1 (ss_1)) Conj Simp2 
       (ss_1) : that .p_1 & .q_1 & .r_1)]

   Assoctest : [(.p_1 : prop), (.q_1 
       : prop), (.r_1 : prop), (ss_1 
       : that (.p_1 & .q_1) & .r_1) => 
       (--- : that .p_1 & .q_1 & .r_1)]

   {move 0}
end Lestrade execution


\end{verbatim}

This text shows that given $(p \wedge q) \wedge r$ we can deduce $p \wedge (q \wedge r)$.  This demonstrates use of the environment extension mechanism:  in the {\tt open...close} block we are able to manipulate {\tt ss} as a constant.

Note that the default grouping of parentheses in nested infix expressions is to the right, so the parentheses are not expressed in the typing of the conclusion.

In the declaration of {\tt Assoctest} you can read the quite complicated expression for the construction which we carried out step by step in the {\tt open...close} block.  Reading this would be made easier by noticing that
Lestrade treats {\tt Conj} as an infix operator!

NOTE:  I would like to turn on full display so one can see what is happening the block.

\end{document}


   
(* quit *)
